\thispagestyle{empty}

\begin{center}
{\large\bfseries Pablo Millán Cubero }\\
\end{center}
\begin{center}
Pablo Millán Cubero\\
\end{center}

%\vspace{0.7cm}

\vspace{0.5cm}
\noindent\textbf{Palabras clave}: \textit{desarrollo de videojuegos, metodologías de desarrollo, godot, gdscript, git, aseprite, software libre, metroidvania}
\vspace{0.7cm}

\noindent\textbf{Resumen}\\

Pese a su importancia y gran impacto económico, la industria del videojuego arrastra problemas de producción que no son tan frecuentes en otros sectores del ocio como el cine y la música. La complejidad software, las necesidades creativas y los grandes equipos multidisciplinares son algunas de las razones por las que es común observar desarrollos en los que la planificación fracasa, se sobrepasa el presupuesto o se implementan juegos de mala calidad. En este trabajo se investigan las particularidades del desarrollo de videojuegos para entender en qué se diferencia de el de software generalista. Con esas conclusiones se propone una metodología y planificación con la que desarrollar un juego de calidad dentro de un plazo fijo. Posteriormente, se realiza una implementación del mismo en el motor Godot.


\cleardoublepage

\begin{center}
	{\large\bfseries Development of a \textit{metroidvania} videogame in Godot}\\
\end{center}
\begin{center}
	Pablo Millán Cubero\\
\end{center}
\vspace{0.5cm}
\noindent\textbf{Keywords}: \textit{game development, godot, gdscript, git, aseprite, open software}, \textit{floss}
\vspace{0.7cm}

\noindent\textbf{Abstract}\\


\cleardoublepage

\thispagestyle{empty}

\noindent\rule[-1ex]{\textwidth}{2pt}\\[4.5ex]

D. \textbf{Pablo García Sánchez}, Profesor(a) del ...

\vspace{0.5cm}

\textbf{Informo:}

\vspace{0.5cm}

Que el presente trabajo, titulado \textit{\textbf{Desarrollo de un videojuego en Godot}},
ha sido realizado bajo mi supervisión por \textbf{Pablo Millán Cubero}, y autorizo la defensa de dicho trabajo ante el tribunal
que corresponda.

\vspace{0.5cm}

Y para que conste, expiden y firman el presente informe en Granada a Junio de 2026.

\vspace{1cm}

\textbf{El/la director(a)/es: }

\vspace{5cm}

\noindent \textbf{(Pablo García Sánchez)}

\chapter*{Agradecimientos}

A todos mis amigos que me ayudaron probando el juego: Pablo, Pablo, Jorge, Simple, Pepe, Antonio, Jaime, Finn, Ale, Oscar, Cof, Carlos, Rocío y Manuel.

Y a Jaime, porque de no ser por sus clases no haría juegos.


